\documentclass[11pt,a4paper]{article}
\usepackage[T1]{fontenc}
\usepackage{lmodern,amsmath,amssymb,amsthm,microtype}
\usepackage[margin=27mm]{geometry}
\usepackage[colorlinks=true,linkcolor=blue,citecolor=blue,urlcolor=blue]{hyperref}
\setlength{\emergencystretch}{2em}
\newtheorem{theorem}{Theorem}[section]
\newtheorem{lemma}[theorem]{Lemma}
\newtheorem{corollary}[theorem]{Corollary}
\newcommand{\Q}{\mathbb Q}\newcommand{\Z}{\mathbb Z}
\DeclareMathOperator{\Res}{Res}\DeclareMathOperator{\Tr}{Tr}
\title{Genus of the actual-trace fibers in a derivative detector}
\author{}\date{17 November 2026}
\begin{document}\maketitle

\begin{abstract}
For the nine-contact derivative detector, fixing a guarded auxiliary
elliptic point gives an actual-trace double cover of the slope line
with at least eighteen simple rational branch points. Its smooth
projective normalization has genus at least eight. The contact
values of the weight and trace-radicand derivatives are respectively
$\pm18$ and $\pm36$. More generally, no geometric rational curve
meets the guarded locus after imposing the elliptic multiplication
relation. These statements exclude rational parametrizations of
that locus; they do not exclude individual rational witnesses or
prove failure of the actual-trace realization criterion.
\end{abstract}

\section{Contact fibers}

Let
\[
                \mathcal S=\{-10,-8,-5,1,2,3,4,6,7\}
\]
and fix a rational number $a>10$. Define
\begin{equation}\label{eq:contact}
 F(a,z)=\prod_{c\in\mathcal S}((a+c)^2-z^2),\qquad
                 u(z)=-18\frac{F(a,z)}{F_a(a,z)},
\end{equation}
where $F_a$ is taken with $z$ fixed. The parameter $a$ is then held
constant when differentiating in $z$.

For each of two slots $\sigma\in\{+,-\}$, fix constants
$C_{1,\sigma},\ldots,C_{6,\sigma}$ in a characteristic-zero field.
Assume the first four are nonzero; the last two may vanish. Put
\begin{equation}\label{eq:detector}
 \begin{gathered}
 D_{i,\sigma}=\frac{C_{i,\sigma}^2}{1+uC_{i,\sigma}^2},\qquad
 W_\sigma=\prod_{i=1}^6\left(1+\frac{D_{i,\sigma}^2}{u}\right),\\
 V=W_+^{-1}+W_-^{-1},\qquad
 w=\frac{u}{(1+u)^2},\qquad
                     \mathcal C=V^2+w(2-V).
 \end{gathered}
\end{equation}
These are rational functions; identities at a contact mean their
regular extensions where such extensions exist.

\begin{lemma}\label{lem:contacts}
The eighteen numbers $z_{c,\varepsilon}=\varepsilon(a+c)$,
with $c\in\mathcal S$ and $\varepsilon\in\{1,-1\}$, are
distinct rational numbers. At each of them,
\[
 F=0,\qquad F_a\ne0,\qquad u=0,\qquad u'=18\varepsilon.
\]
Moreover $\mathcal C$ is regular there, with
\[
                       \mathcal C=0,\qquad
                         \mathcal C'=36\varepsilon.
\]
\end{lemma}
\begin{proof}
Every $a+c$ is positive and the nine values are distinct. Their
positives and negatives therefore give eighteen distinct nonzero
rational points. Set
\[
 P_c(a)=\prod_{d\in\mathcal S\setminus\{c\}}
                ((a+d)^2-(a+c)^2)
       =\prod_{d\ne c}(d-c)(2a+c+d).
\]
Each factor is nonzero: $d\ne c$ and $2a+c+d>0$ because $a>10$.
At $z=z_{c,\varepsilon}$, direct differentiation gives
\[
        F_a=2(a+c)P_c(a),\qquad
        F_z=-2\varepsilon(a+c)P_c(a).
\]
Thus $u$ has a simple zero and $u'=-18F_z/F_a=18\varepsilon$.

In a slot with $r_\sigma$ nonzero tests, $4\le r_\sigma\le6$,
each nonzero test contributes
\[
 \left(1+\frac{D_{i,\sigma}^2}{u}\right)^{-1}
 =\frac{u(1+uC_{i,\sigma}^2)^2}
        {u(1+uC_{i,\sigma}^2)^2+C_{i,\sigma}^4}.
\]
This has a simple zero at $u=0$; a zero test contributes one.
Consequently $W_\sigma^{-1}=O(u^{r_\sigma})$ and
$V=O(u^4)$. Possible cancellation between the slots only raises
the latter order. Therefore
\begin{equation}\label{eq:expansion}
                         \mathcal C=2u+O(u^2).
\end{equation}
Composing with \eqref{eq:contact} gives
$\mathcal C'=2u'=36\varepsilon$.
\end{proof}

\begin{theorem}\label{thm:genus}
Suppose the two slots are conjugate over a quadratic number field,
so $\mathcal C(z)\in\Q(z)$. The double cover
\[
                              Y^2=\mathcal C(z)
\]
has a geometrically integral smooth projective normalization of
genus at least eight. Its eighteen contact points are simple rational
branch points, and their unique points on the normalization are
rational. The same geometric genus bound holds after fixing arbitrary
slot constants over any algebraically closed field of characteristic
zero subject to the four-nonzero-tests hypothesis.
\end{theorem}
\begin{proof}
Lemma~\ref{lem:contacts} gives eighteen simple zeros of
$\mathcal C$. Hence its square class is nontrivial even over an
algebraic closure, and the associated function-field extension is
geometrically quadratic. At a simple zero the local equation has
nonzero derivative in $z$, so its single point with $Y=0$ is smooth
and ramified of index two. Over a rational contact its residue field
is $\Q$.

Let $r$ be the total number of geometric branch points of the
normalized double cover, including all additional zeros and poles
of odd order. Riemann--Hurwitz gives
\[
                         2g-2=-4+r,\qquad r\ge18.
\]
Thus $g\ge8$. The same proof uses only characteristic zero and the
stated nonzero conditions on the constants.
\end{proof}

The original source requires $u\ne0$. Accordingly all eighteen
known rational ramification points are outside its guarded locus.
The theorem does not assert the existence or absence of any other
rational point.

\section{The elliptic multiplication locus}

Fix the curve and point of \cite{Derivative,Criterion}:
\[
 K=\Q(\sqrt2),\quad E:y^2=x^3-(11-6\sqrt2)x,\quad
 P=\left(\frac{57-28\sqrt2}{8},\frac{-660+417\sqrt2}{32}\right).
\]
Fix an odd $B\ge3$ and a nonzero integer $n$. Set
\[
 A=[n]P,\qquad U=[n]H,\qquad
 T=\Tr x(A),\qquad a=((n^2+B)^4+T^4)^4.
\]
The point $A$ and the rational $a$ are fixed while $H$ and $z$ vary.
In particular $a\ge B^{16}>10$.
Use $f=1/x$, $\zeta=-x/y$ and the six actual tests
\[
 \begin{gathered}
 C_1=f(H),\quad C_2=f(U),\quad
 C_3=\overline{f(H-P)},\quad C_4=\overline{f(U-A)},\\
 C_5=\frac{\zeta(U)}{\zeta(H)}-
          \overline{\frac{\zeta(U-A)}{\zeta(H-P)}},\qquad
 C_6=\frac{\zeta(U)}{\zeta(H)}-n.
 \end{gathered}
\]
All original point, affine-difference and quotient guards are retained.
They make every test finite. Each of the first four is a reciprocal
of the finite nonzero abscissa of an affine point, so is nonzero in
each geometric slot. This remains true when the last two tests vanish.
Theorem~\ref{thm:genus} therefore applies to every guarded fixed $H$.

For a precise geometric statement, let $R=\Res_{K/\Q}E$, an abelian
surface, and let $R^\circ$ be the open subset on which all these
point and quotient expressions are defined and the first four tests
are nonzero in both slots. The multiplication relation uses the
morphism $[n]:R\to R$; it introduces no independent $U$ coordinate.
The rational functions \eqref{eq:contact}--\eqref{eq:detector} are
then defined over $R^\circ\times\mathbf P^1_z$.
Let $\mathcal X^\circ_{B,n}$ be the locus $Y^2=\mathcal C$
on which all retained source denominators and nonzero guards hold.
It is defined by the rational equation on this open set, without
extra components from clearing a vanishing denominator.

\begin{corollary}\label{cor:rational}
There is no nonconstant rational map
$\mathbf P^1_{\overline\Q}\dashrightarrow\mathcal X^\circ_{B,n}$.
Thus no geometric rational curve has a nonempty open subset in this
guarded locus. The same assertion holds for any finite point-chart
cover of the guarded multiplication locus.
\end{corollary}
\begin{proof}
Project such a map to $R$. Since $R$ is proper, the rational map
from the smooth curve $\mathbf P^1$ extends to a morphism. It is
constant: after embedding in $\mathbb C$, translation-invariant
holomorphic one-forms span every cotangent space of the abelian
surface, while $\mathbf P^1$ has no holomorphic one-forms. The
morphism has zero differential and is constant in characteristic zero.

Thus the proposed curve lies over a fixed $H\in R^\circ$.
If its $z$ coordinate varies, it induces a nonconstant map from
$\mathbf P^1$ to the smooth projective normalization of that fiber,
whose genus is at least eight by Theorem~\ref{thm:genus}.
Riemann--Hurwitz rules out this map. If $z$ is constant, the quadratic
equation gives at most two values of $Y$, so the image is a point.
A finite point-chart cover cannot have a curve mapping to a point
of this locus. The same conclusion follows for such a cover.
\end{proof}

This excludes rational parametrizations of the guarded multiplication
locus, including ones that attempt to vary $H$ rationally along with
the existing slope. It neither excludes isolated rational witnesses
nor proves that there are only finitely many witnesses as $H$ varies.
In particular it does not disprove $\mathrm{TC}_B(n)$ or its eventual
positive-input form in \cite{Criterion}. These statements concern the displayed derivative and trace test;
they assert no such obstruction for a different detector.

\section{Exact verification}

The checker
\nolinkurl{negative_results/computation/derivative_trace_genus_v65/check.py}
uses exact polynomial arithmetic to verify all eighteen contact
identities, nonzero factor decompositions, and the values
$u'=18\varepsilon$, $\mathcal C'=36\varepsilon$.
It also checks the reciprocal-factor identity and the first-order
trace expansion. The normalization, ramification and rational-curve
arguments are proved above, not inferred from numerical samples.

\begin{thebibliography}{9}\raggedright
\bibitem{Derivative}\emph{An eleven-variable construction under rational
$abc$}, repository manuscript, 17 November 2026, version 63,
\nolinkurl{papers/conditional_few_variables/derivative_detector11_20261117_v63.tex},
published commit
\nolinkurl{f70057ffce6b8318a79a7e3648779067fd98c2e3}.
\bibitem{Criterion}\emph{One actual-trace square in a derivative detector},
repository manuscript, 17 November 2026, version 64,
\nolinkurl{papers/pole_parameters/single_trace_derivative_criterion_20261117_v64.tex},
published commit
\nolinkurl{aa7013ce6129af0db621b89d5e44deaa07590b35}.
\end{thebibliography}
\end{document}
